%% Elsevier CAS double-column manuscript draft
%% Project: PI-MSCL battery SOH estimation

\documentclass[a4paper,fleqn]{cas-dc}

\usepackage[numbers]{natbib}
\usepackage{amsmath,amssymb}
\usepackage{booktabs}
\usepackage{multirow}
\usepackage{graphicx}
\usepackage{siunitx}

\begin{document}
\let\WriteBookmarks\relax
\def\floatpagepagefraction{1}
\def\textpagefraction{.001}

\shorttitle{PI-MSCL for battery SOH estimation}
\shortauthors{L. Chang et~al.}

\title[mode=title]{Physics-Informed Multi-Scale CNN-LSTM for Robust Lithium-Ion Battery State-of-Health Estimation}

\author[1]{Chang Liu}
\cormark[1]
\ead{liuchang2262@gmail.com}
\credit{Conceptualization, Methodology, Software, Validation, Writing - original draft}

\affiliation[1]{organization={},
                addressline={},
                city={},
                postcode={},
                country={}}

\cortext[cor1]{Corresponding author}

\begin{abstract}
Accurate state-of-health (SOH) estimation is essential for lithium-ion battery management, but capacity-derived SOH labels are substantially more costly to obtain than operating features. This study formulates a controlled sparse-label benchmark on the fully labeled HUST dataset by randomly retaining a fixed proportion of cycle-level targets within every training cell while preserving complete feature trajectories and fully labeled unseen-cell evaluation. A physics-informed multi-scale CNN-LSTM, termed PI-MSCL, combines three temporal convolutional receptive fields, recurrent cross-cycle modeling, masked data fitting, and a soft monotonic trajectory constraint. A unified $2\times2$ experiment separates the effects of multi-scale representation and physical consistency over four label-retention ratios, and evaluates both pointwise error and trajectory plausibility. Numerical findings will be inserted only after the leakage-free multi-seed experiment matrix is complete.
\end{abstract}

\begin{highlights}
\item A controlled within-battery sparse-label protocol is defined on fully labeled HUST data.
\item Label masking preserves operating features and cycle order at unlabelled targets.
\item A shared $2\times2$ design separates multi-scale and physical-constraint effects.
\item Pointwise accuracy is evaluated together with trajectory-level plausibility.
\end{highlights}

\begin{keywords}
Lithium-ion battery \sep State of health \sep Partial lifecycle supervision \sep Physics-informed learning \sep Multi-scale CNN-LSTM \sep Battery management system
\end{keywords}

\maketitle

\section{Introduction}

Lithium-ion batteries have become core energy-storage components in electric vehicles, grid storage, and portable electronics, making reliable battery health monitoring a prerequisite for safe and economical operation. State of health (SOH), commonly expressed through the relative loss of available capacity or power capability, is one of the most important indicators used by battery management systems (BMSs) to support safety warning, lifetime management, and maintenance decisions \cite{berecibar2016critical}. Recent reviews of second-life lithium-ion battery applications also emphasize that SOH estimation remains a core bottleneck for safe reuse and operation \cite{vignesh2024state}. Degradation diagnostics further show that aging changes are governed by coupled electrochemical and mechanical processes \cite{birkl2017degradation}. These processes, including loss of active material, solid-electrolyte interphase growth, lithium inventory loss, and impedance increase, change the voltage, current, temperature, and charge-discharge response observed during operation \cite{che2023health}. Multiscale imaging and computational modeling studies further show that degradation mechanisms can evolve across coupled material and electrode scales \cite{zhang2023coupling}, while knee-point analyses demonstrate that the rate of capacity loss can change markedly within a lifecycle \cite{attia2022knees}. These measurable signals create an opportunity for data-driven SOH estimation, but they also expose a central difficulty: the observable operating trajectory is continuous, whereas reliable capacity calibration is often intermittent, costly, and unavailable over the full lifecycle.

Existing SOH estimation studies have developed along two broad lines. Model-based and diagnostic methods, including equivalent-circuit analysis, incremental capacity analysis, and physics-based degradation interpretation, offer useful physical insight but often rely on controlled testing conditions, careful feature extraction, or parameter identification procedures that are difficult to maintain in long-term online deployment \cite{dubarry2022best}. Data-driven methods reduce this burden by learning the mapping from measurable charge-discharge signals to SOH. Early-cycle degradation trajectories have been used for data-driven lifetime prediction \cite{severson2019data}, while machine-learning pipelines and general data-driven battery models have demonstrated the value of operating-feature representations \cite{roman2021machine}. Data-driven models have also been used to jointly predict battery state of charge and health from measurable signals \cite{ng2020predicting}. Partial charging curves further show that useful health information can be extracted from practical charge-discharge segments \cite{lyu2021partial}. Robust CNN and random-forest combinations indicate that feature learning and ensemble regression can improve SOH estimation under noisy operating data \cite{yang2022robust}. A recent review further calls for unified machine-learning characterization of battery degradation across multiple levels \cite{li2022towards}. More recent deep learning approaches, including recurrent neural networks, convolutional networks, CNN-LSTM hybrids, attention-based models, and transfer-learning frameworks, further improve representation learning for nonlinear degradation patterns and cross-cell variability \cite{hochreiter1997long}. Online LSTM-based estimation is a representative example of recurrent health modeling from charging-process data \cite{liu2023online}. Attention-based CNN-BiLSTM models further combine local curve representation with temporal dependency learning for SOH and RUL estimation \cite{zhu2022attention}. Transformer-LSTM architectures have also been explored for fast-charging SOH estimation \cite{fan2023online}. Temperature-prediction GRU models provide another route for embedding thermal response information into SOH estimation \cite{chen2022state}. Hybrid attention and deep-learning methods further indicate the value of adaptive temporal feature weighting \cite{zhao2023state}. Deep learning without additional degradation experiments has been investigated as a way to reduce experimental burden in battery health estimation \cite{lu2023deep}. Deep transfer learning has also been explored to improve personalized battery health prediction across cells \cite{ma2022real}. Cycle-synchronization transfer learning addresses cross-cell temporal misalignment in SOH estimation \cite{zhou2022transfer}. Graph neural networks based on partial discharging curves further suggest that relational battery representations can support health estimation when complete curves are unavailable \cite{zhou2024graph}. These methods have substantially advanced battery SOH estimation, especially when sufficiently dense labels are available.

However, most supervised data-driven models are developed with dense target coverage: each training cycle or sampled window is paired with a capacity-derived SOH value. This assumption is convenient for benchmark construction but expensive to reproduce, because reliable SOH targets generally depend on reference capacity measurements rather than direct sensing. Recent label-efficient studies confirm that this bottleneck is substantive, but they instantiate label scarcity in different ways. Few-shot learning reduces the number of training examples available to a Bayesian estimator \cite{zhang2023fewshot}, whereas early-life labeled observations can be expanded through feature-informed data augmentation \cite{han2024limited}. Self-supervised learning with weak labels uses auxiliary degradation signals before fine-tuning on a smaller annotated set \cite{wang2024weaklabels}, and weakly supervised pre-training has recently been extended to extremely small sets of real SOH labels \cite{qian2026minimal}. Semi-supervised transfer learning \cite{li2020state}, unlabeled charging data \cite{lin2023improving}, maritime battery monitoring \cite{salucci2023novel}, and selective pseudo-labeling with auxiliary tasks \cite{li2026auxiliary} provide further evidence that unannotated cycling data can support battery health estimation. These studies establish label efficiency as an important objective, but their supervision units and auxiliary information differ: scarcity may occur across cells, early-life samples, isolated curve segments, or weak and pseudo labels. The present study instead isolates a cycle-level SOH-label budget within every training cell while retaining the complete sequence of operating features and cycle order.

Physics-informed learning provides a complementary mechanism for this setting because trajectory constraints do not require a true SOH target at every cycle. In the broader machine-learning literature, physics-informed neural networks and physics-informed machine learning incorporate governing relations or domain constraints into the learning objective to improve data efficiency and the admissibility of learned solutions \cite{raissi2019physics,karniadakis2021physics}. In battery applications, physics-guided models have been developed for SOH prognostics from partial charging segments \cite{kohtz2022physics}, stable degradation modeling and prognosis \cite{wang2024physics}, and neural-network estimation with explicit degradation constraints \cite{ye2024method}. Physics-informed prognostics and health-management studies further demonstrate the value of embedding battery knowledge into learned health representations \cite{wen2023physics}. Recent work has extended physics-aware estimation to partial charging profiles \cite{gu2026real}, varying operating conditions \cite{lin2025physics}, and different cell chemistries \cite{lin2025lightweight}. For sparse cycle-level supervision, physical consistency can therefore play a more specific role than an auxiliary penalty evaluated only at labeled samples: it can provide structural weak supervision along the retained feature trajectory. Because short-term capacity recovery and measurement fluctuation can produce local SOH increases, this prior should be implemented softly rather than as a hard monotonic rule.

In this paper, we propose a physics-informed multi-scale CNN-LSTM framework, termed PI-MSCL, for lithium-ion battery SOH estimation under controlled sparse-label supervision. We use the completely labeled HUST dataset to construct a transparent test bed: batteries are first separated into training, validation, and test sets, after which a fixed proportion of cycle-level SOH targets is randomly retained within each training battery. Input features and cycle order remain complete, and validation and test batteries remain fully labeled. This protocol is an artificial intervention for measuring robustness to a reduced label budget; it is not presented as a reproduction of a naturally missing field-data process. PI-MSCL uses parallel one-dimensional convolutional branches with different receptive fields to extract short-, medium-, and long-range degradation patterns, followed by an LSTM to model cross-cycle health evolution. A masked supervised loss uses only retained labels, while a soft monotonic consistency term penalizes excessive upward deviations along the complete predicted trajectory after an early-cycle threshold. The two objectives thus combine calibrated observations with label-independent structural information.

The main contributions of this study are as follows:
\begin{itemize}
  \item We define a reproducible, battery-stratified cycle-level masking protocol on HUST that varies the training-label budget while preserving complete operating trajectories and fully labeled cross-cell evaluation.
  \item We develop PI-MSCL, a compact multi-scale CNN-LSTM that combines three temporal receptive fields with recurrent cross-cycle modeling and soft monotonic structural supervision.
  \item We use a controlled two-factor comparison to separate the contribution of multi-scale representation from that of physical consistency under four label-retention ratios.
  \item We evaluate not only pointwise SOH error but also trajectory-level monotonicity and upward-deviation behavior, allowing accuracy gains and physical plausibility to be interpreted without conflation.
\end{itemize}

The remainder of this paper is organized as follows. Section~\ref{sec:methodology} presents PI-MSCL and its learning objectives. Section~\ref{sec:experiments} defines the HUST dataset, sparse-label construction, comparison design, and evaluation metrics. Section~\ref{sec:results} reports the full-supervision benchmark, label-budget results, trajectory analysis, and ablation study. Section~\ref{sec:conclusion} summarizes the findings and their evidence boundaries.

\section{Methodology}
\label{sec:methodology}

\subsection{Problem formulation and learning interface}
\label{subsec:problem_formulation}

The target task is to estimate the SOH of an unseen lithium-ion cell from operating features extracted from its charging process. For a battery cell indexed by $i$, let $\mathbf{x}_{i,t}\in\mathbb{R}^{F}$ denote the feature vector extracted at cycle or window index $t$, where $F$ is the feature dimension. A model input is a fixed-length feature sequence
\begin{equation}
\mathbf{X}_{i,t}=\left[\mathbf{x}_{i,t-T+1},\mathbf{x}_{i,t-T+2},\ldots,\mathbf{x}_{i,t}\right]\in\mathbb{R}^{T\times F},
\label{eq:input_sequence}
\end{equation}
where $T$ is the sequence length and $F$ is the input-feature dimension. The corresponding target is the normalized SOH value $y_{i,t}\in[0,1]$ at the final cycle of the window. Input windows are constructed independently within each battery, so no sequence can cross a cell boundary. Capacity and SOH are reserved for target construction and are never included in $\mathbf{X}_{i,t}$.

The model receives a binary supervision indicator $m_{i,t}\in\{0,1\}$ together with each target position. An indicator of one permits direct data fitting, whereas an indicator of zero disables only the target-dependent term. The feature window, battery identity, and cycle index remain available in both cases. For a mini-batch $\mathcal{B}$, the masked supervised loss is
\begin{equation}
\mathcal{L}_{\mathrm{sup}}=
\begin{cases}
\dfrac{\sum_{(i,t)\in\mathcal{B}} m_{i,t}\left(\hat{y}_{i,t}-y_{i,t}\right)^2}
{\sum_{(i,t)\in\mathcal{B}} m_{i,t}}, & \sum_{(i,t)\in\mathcal{B}}m_{i,t}>0,\\[6pt]
0, & \text{otherwise},
\end{cases}
\label{eq:masked_mse}
\end{equation}
where $\hat{y}_{i,t}$ is the predicted SOH. A mini-batch without a labeled target contributes zero data loss but can still contribute physical loss. When all target labels are available, Eq.~\eqref{eq:masked_mse} reduces to the standard mean squared error. The experimental procedure used to construct $m_{i,t}$ and control the label budget is specified in Section~\ref{subsec:split_protocol}, rather than being treated as part of the model architecture.

\subsection{Overview of PI-MSCL}
\label{subsec:overview}

PI-MSCL is designed to connect three requirements of robust battery SOH estimation: extracting degradation features at multiple temporal scales, modeling cross-cycle health evolution, and constraining the predicted trajectory with physically admissible degradation behavior. As illustrated conceptually in Fig.~\ref{fig:pimscl_framework}, the framework consists of four components: a multi-scale one-dimensional convolutional feature extractor, an LSTM temporal modeling module, a bounded SOH regression head, and a physics-informed training objective.

\begin{figure*}[t]
  \centering
  \includegraphics[width=\textwidth]{../../docs/PI-MSCL_architecture_v3_preview.pdf}
  \caption{PI-MSCL under controlled sparse-label supervision.}
  \label{fig:pimscl_framework}
\end{figure*}

Figure~\ref{fig:pimscl_framework} follows the evidence path of the method. Panel (a) distinguishes sparse SOH calibration points from the complete operating-feature trajectory. Panel (b) expands the learnable mapping from the input tensor through three convolutional scales, feature fusion, cross-cycle LSTM modeling, and the bounded regression head. Panel (c) separates the label-supported data path from the trajectory-structure path. Thus, a cycle without a retained SOH target can still affect the learned representation as an input and can participate in the soft monotonic constraint through its prediction.

\subsection{Multi-scale convolutional feature extraction}
\label{subsec:multiscale}

Battery degradation signals contain information at different temporal scales. Short-range variations reflect local charge-response fluctuations and measurement dynamics between neighboring cycles. Medium-range patterns describe transitional changes over several cycles, such as shifts in charge duration or voltage statistics. Long-range patterns reflect the slower evolution of capacity fade and resistance-related behavior across a wider lifecycle window. A single convolutional kernel can only emphasize one receptive-field scale, which may be insufficient when the same feature sequence contains both local fluctuations and slowly varying degradation trends.

To address this issue, PI-MSCL uses a three-branch multi-scale one-dimensional CNN module. Each branch receives the same input sequence but uses a different kernel size:
\begin{equation}
\mathbf{H}^{(s)}_{i,t}=\phi_{s}\left(\mathbf{X}_{i,t};k_s\right),
\quad k_s\in\{3,7,15\},
\label{eq:multiscale_branch}
\end{equation}
where $\phi_s(\cdot)$ denotes a branch containing two repeated Conv1D--batch-normalization--ReLU--max-pooling blocks. Same padding is used in the convolutions, and each pooling layer has kernel size and stride two. The three kernel sizes represent short-, medium-, and long-range receptive fields, respectively, while both convolutional layers in every branch contain 64 channels. For $T=40$, the two pooling operations reduce each aligned branch sequence to 10 temporal positions. The three outputs are concatenated along the channel dimension and fused by a $1\times1$ convolution:
\begin{equation}
\mathbf{H}_{i,t}=\phi_{\mathrm{fuse}}\left(
\left[\mathbf{H}^{(3)}_{i,t};\mathbf{H}^{(7)}_{i,t};\mathbf{H}^{(15)}_{i,t}\right]
\right),
\label{eq:multiscale_fusion}
\end{equation}
where $[\cdot;\cdot;\cdot]$ denotes channel-wise concatenation and $\phi_{\mathrm{fuse}}(\cdot)$ maps the resulting 192 channels to 128 channels using Conv1D, batch normalization, and ReLU activation. This design changes the receptive-field structure without changing the downstream LSTM used by the single- and multi-scale variants.

\subsection{LSTM-based temporal degradation modeling}
\label{subsec:lstm}

After multi-scale feature fusion, the output sequence $\mathbf{H}_{i,t}$ is fed into an LSTM module to model the temporal evolution of battery health. LSTM networks are suitable for sequential modeling because their gated state updates can preserve relevant historical information while filtering transient noise \cite{hochreiter1997long}. In PI-MSCL, a two-layer LSTM with hidden dimension 64 is used:
\begin{equation}
\mathbf{Z}_{i,t}=\mathrm{LSTM}\left(\mathbf{H}_{i,t}\right),
\label{eq:lstm_sequence}
\end{equation}
where $\mathbf{Z}_{i,t}$ is the sequence of hidden states. The last hidden state $\mathbf{z}_{i,t}^{\mathrm{last}}$ is treated as the degradation representation of the input window. It is passed through a fully connected regression head,
\begin{equation}
\hat{y}_{i,t}=\sigma\left(f_{\mathrm{fc}}\left(\mathbf{z}_{i,t}^{\mathrm{last}}\right)\right),
\label{eq:soh_regression}
\end{equation}
where $f_{\mathrm{fc}}(\cdot)$ comprises a 64-unit fully connected layer, ReLU activation, dropout with probability 0.4, and a scalar output layer; $\sigma(\cdot)$ is the sigmoid function. The sigmoid mapping bounds the predicted SOH to $[0,1]$, consistent with normalized capacity-based SOH, and makes a separate boundary penalty unnecessary in the main configuration.

\subsection{Physics-informed training objective}
\label{subsec:objective}

The supervised loss in Eq.~\eqref{eq:masked_mse} learns the numerical mapping between operating features and SOH when calibrated labels are available. However, when labels are sparse, the data-fitting term provides no direct gradient for unlabeled cycles. PI-MSCL therefore adds a soft monotonic physical consistency constraint to the predicted SOH trajectory. The constraint is based on the macroscopic degradation tendency that, after the early activation and measurement-fluctuation stage, battery SOH should generally follow a non-increasing trend.

For a mini-batch, define the set of valid ordered pairs
\begin{equation}
\mathcal{P}=\left\{(a,b): i_a=i_b,\;0<c_b-c_a\leq K,\;c_a\geq c_{\min}\right\},
\label{eq:pair_set}
\end{equation}
where $i_a=i_b$ means that the two samples belong to the same battery, $c_a$ and $c_b$ are their cycle indices, $K$ is the maximum cycle gap considered in the pairwise constraint, and $c_{\min}$ is the delayed activation threshold. The condition $c_a\geq c_{\min}$ prevents early-cycle capacity recovery or activation behavior from being over-penalized. The soft monotonic loss is
\begin{equation}
\mathcal{L}_{\mathrm{mono}}=
\frac{1}{|\mathcal{P}|}
\sum_{(a,b)\in\mathcal{P}}
\omega_{a,b}\,
\mathrm{ReLU}\left(\hat{y}_{i_b,c_b}-\hat{y}_{i_a,c_a}-\epsilon\right),
\label{eq:mono_loss}
\end{equation}
where $\epsilon$ is the tolerance for small local upward fluctuations and $\omega_{a,b}$ is a temporal decay weight. The pair set is constructed within each mini-batch using battery and true cycle metadata; supervision indicators do not enter Eq.~\eqref{eq:mono_loss}. An exponential decay is used:
\begin{equation}
\omega_{a,b}=\exp\left[-\alpha\left(c_b-c_a\right)\right],
\label{eq:temporal_decay}
\end{equation}
which gives stronger weight to nearby cycle pairs than to distant pairs. The main configuration uses $\epsilon=0.005$, $c_{\min}=300$, $K=40$, and $\alpha=0.2$.

The full training objective is
\begin{equation}
\mathcal{L}=\mathcal{L}_{\mathrm{sup}}
+\lambda_{\mathrm{mono}}\mathcal{L}_{\mathrm{mono}},
\label{eq:total_loss}
\end{equation}
where $\lambda_{\mathrm{mono}}=0.3$ in the main PI-MSCL configuration. Although the training library supports additional boundary and second-order smoothness terms, both are disabled in the main experiment. The reported method therefore combines the masked data loss with one explicit trajectory prior, making the architectural and physical contributions identifiable in the factorial comparison.
\section{Experimental Setup}
\label{sec:experiments}

\subsection{Datasets and SOH definition}
\label{subsec:datasets}

The study uses the public HUST lithium-ion battery degradation dataset \cite{Yuan2022HUSTDataset}. It contains full-lifecycle records from 77 LFP/graphite cells with a nominal capacity of 1.1 Ah and nominal voltage of 3.3 V. The cells were tested at 30~$^{\circ}$C using a common constant-current/constant-voltage charging procedure and multiple discharge profiles. The resulting collection contains substantial variation in lifetime, fading rate, and local capacity fluctuation across cells, allowing a battery-level split to test generalization to unseen cells. Using a public dataset also supports reproducible degradation analysis and is consistent with calls for standardized battery-data infrastructure \cite{ward2022principles}.

SOH was defined as normalized capacity relative to the initial capacity of each cell:
\begin{equation}
\mathrm{SOH}_{i,t}=\frac{Q_{i,t}}{Q_{i,0}},
\label{eq:soh_definition}
\end{equation}
where $Q_{i,t}$ is the measured capacity of cell $i$ at cycle $t$ and $Q_{i,0}$ is the initial measured capacity. This definition makes the target dimensionless and aligns the prediction range with the sigmoid-bounded output of PI-MSCL.

HUST provides capacity measurements for all recorded cycles. Sparse supervision is therefore introduced only after the battery split as a controlled intervention on the training targets. The original measurements remain unchanged, and no HUST sample is presented as naturally unlabeled field data.

\subsection{Feature construction and preprocessing}
\label{subsec:features_preprocessing}

% TBC-008: the Chinese draft states 14 features, whereas the audited loader and
% local HUST CSV files use the 16 columns reported below. Retain this explicit
% implementation-aligned version until the author confirms the final feature set.
The audited implementation reads the 16 charging-cycle features summarized in Table~\ref{tab:features}. They describe voltage and current statistics, CC/CV charge throughput and duration, local signal slopes, and entropy-based complexity. The discharge-capacity column used to construct SOH is excluded from the input, preventing direct target leakage.

\begin{table*}[t]
\caption{Charging-cycle features used for SOH estimation.}
\label{tab:features}
\centering
\begin{tabular}{lll}
\toprule
Group & Feature & Description \\
\midrule
Voltage statistics & voltage mean & Mean voltage during the CC stage \\
Voltage statistics & voltage std & Voltage fluctuation during the CC stage \\
Voltage statistics & voltage kurtosis & Sharpness of the voltage distribution \\
Voltage statistics & voltage skewness & Asymmetry of the voltage distribution \\
CC-stage dynamics & CC Q & Charged capacity during the CC stage \\
CC-stage dynamics & CC charge time & Duration of the CC charging stage \\
Voltage dynamics & voltage slope & Local slope of the charging-voltage signal \\
Voltage complexity & voltage entropy & Entropy-based complexity of the voltage signal \\
Current statistics & current mean & Mean current during the CV stage \\
Current statistics & current std & Current fluctuation during the CV stage \\
Current statistics & current kurtosis & Sharpness of the current distribution \\
Current statistics & current skewness & Asymmetry of the current distribution \\
CV-stage dynamics & CV Q & Charged capacity during the CV stage \\
CV-stage dynamics & CV charge time & Duration of the CV charging stage \\
Current dynamics & current slope & Local slope of the charging-current signal \\
Current complexity & current entropy & Entropy-based complexity of the current signal \\
\bottomrule
\end{tabular}
\end{table*}

No three-sigma row deletion or synthetic data-degradation transform is enabled in the main experiment, so lifecycle continuity is preserved. After the cells are assigned to training, validation, and test sets, a \texttt{StandardScaler} is fitted once using only the concatenated training-cell features. The fitted transformation is then applied unchanged to validation and test cells. In particular, no scaler is fitted independently on an unseen cell. Fixed-length windows of $T=40$ are subsequently generated within each battery, and the SOH at the final position is used as the many-to-one target.

\subsection{Battery-level split and controlled sparse-label protocol}
\label{subsec:split_protocol}

All experiments use a battery-level split rather than a random cycle-level data split. Integer allocation of the 77 HUST cells at a 60/20/20 ratio gives 46 training cells, 15 validation cells, and 16 test cells. All cycles from a given cell remain in one subset. Test cells are therefore unseen during optimization, and the reported errors quantify cross-cell generalization rather than interpolation between cycles of a previously observed cell.

After the split, the label-retention ratio $r$ controls the available SOH targets only within the training cells. For each training battery $i$, a seeded random sampler without replacement selects
\begin{equation}
  n_i^{\mathrm{lab}}=\max\left(1,\operatorname{round}(rN_i)\right)
\end{equation}
of its $N_i$ cycle-level targets and assigns $m_{i,t}=1$ to the selected indices. The remaining target indicators are zero. The experiment uses $r\in\{1.0,0.7,0.5,0.3\}$; $r=1.0$ is the full-supervision reference. Feature rows, cycle order, and windows are never removed by this operation. Validation and test cells remain fully labeled for model selection and final evaluation. Thus, the protocol measures sensitivity to a controlled within-battery label budget; it does not model a continuous missing interval, an early-life prefix, or a periodic field-calibration schedule.

Battery-split, label-mask, and network-initialization seeds are recorded separately for every run. The experiment runner implements both repeated independent battery partitions and a fixed-partition sensitivity protocol, but their estimates are not pooled because they answer different questions. The final result tables will state the selected repetition policy and report the arithmetic mean and sample standard deviation across its completed runs.

\subsection{Compared models and ablation design}
\label{subsec:baselines}

The primary comparison is a controlled $2\times2$ factorial design implemented by one shared model class. The architecture factor switches between a single Conv1D stream with kernel size 7 and the proposed parallel streams with kernel sizes 3, 7, and 15. The physics factor switches the soft monotonic loss off or on. This produces single-scale CNN-LSTM, PI-CNN-LSTM, multi-scale CNN-LSTM, and PI-MSCL. All four variants use identical battery splits, input windows, preprocessing, LSTM and regression heads, optimizer settings, label masks, and stopping rules for a given run.

This design produced a two-dimensional ablation matrix: single-scale versus multi-scale representation on the architecture axis, and no-physics versus soft monotonic physical consistency on the constraint axis. The matrix allowed the experiments to distinguish whether performance gains came primarily from richer degradation representation, from the physical constraint, or from their combination.

\begin{table}[t]
\caption{Model variants used in the main comparisons and ablation analysis.}
\label{tab:model_variants}
\centering
\begin{tabular}{lll}
\toprule
Model & Multi-scale CNN & Soft monotonic constraint \\
\midrule
CNN-LSTM & No & No \\
PI-CNN-LSTM & No & Yes \\
MS-CNN-LSTM & Yes & No \\
PI-MSCL & Yes & Yes \\
\bottomrule
\end{tabular}
\end{table}

\subsection{Training configuration}
\label{subsec:training_config}

All variants use a sequence length of 40. PI-MSCL uses three convolutional branches with kernel sizes 3, 7, and 15; each branch contains two 64-channel layers, and the concatenated output is fused to 128 channels. The common temporal head is a two-layer LSTM with hidden dimension 64, followed by a 64-unit fully connected layer, ReLU, dropout of 0.4, a scalar layer, and sigmoid output. PI-MSCL contains 241,281 trainable parameters for the currently audited 16-feature input. Its physical settings are $\lambda_{\mathrm{mono}}=0.3$, $\epsilon=0.005$, $c_{\min}=300$, $K=40$, and $\alpha=0.2$.

Models are optimized with Adam for at most 200 epochs using a batch size of 1024. The learning rate increases linearly from $2\times10^{-3}$ to $5\times10^{-3}$ over 20 warmup epochs and then follows cosine decay to $10^{-4}$. Early stopping monitors validation MAE with patience 20 and a minimum improvement of $10^{-5}$; the checkpoint with the lowest validation MAE is restored for testing. Every formal run archives the complete configuration, three seed roles, battery lists, checkpoint, predictions, targets, battery identifiers, and true cycle indices. Main tables report mean $\pm$ sample standard deviation; a one-epoch smoke run is marked in metadata and is automatically rejected by the paper-figure script.

\subsection{Evaluation metrics}
\label{subsec:metrics}

SOH estimation performance was evaluated using mean absolute error (MAE), root mean squared error (RMSE), and the coefficient of determination ($R^2$):
\begin{align}
\mathrm{MAE} &= \frac{1}{N}\sum_{n=1}^{N}\left|\hat{y}_{n}-y_{n}\right|, \\
\mathrm{RMSE} &= \sqrt{\frac{1}{N}\sum_{n=1}^{N}\left(\hat{y}_{n}-y_{n}\right)^2}, \\
R^2 &= 1-\frac{\sum_{n=1}^{N}\left(y_n-\hat{y}_n\right)^2}{\sum_{n=1}^{N}\left(y_n-\bar{y}\right)^2}.
\end{align}
MAE measures average prediction deviation, RMSE emphasizes large errors, and $R^2$ measures explained variance. All three are computed over every window of the fully labeled test cells. Trajectory metrics are then computed independently within each test battery after sorting predictions by the saved true cycle index. Let $\Delta\hat{y}_{i,t}=\hat{y}_{i,t+1}-\hat{y}_{i,t}$ for adjacent available predictions whose earlier cycle is at least $c_{\min}$. The monotonicity-violation rate and cumulative upward excess are
\begin{align}
V_{\mathrm{mono}} &= 100\,\frac{\sum_{i,t}\mathbb{I}\!\left(\Delta\hat{y}_{i,t}>\epsilon\right)}{N_{\mathrm{pair}}},\\
U_i &= \sum_t \max\left(\Delta\hat{y}_{i,t}-\epsilon,0\right).
\end{align}
The analysis reports $V_{\mathrm{mono}}$, the mean nonzero violation excess, the mean of $U_i$ across test cells, and the mean absolute second difference as a trajectory-roughness descriptor. The evaluation uses the same $\epsilon=0.005$ and $c_{\min}=300$ as training, while only adjacent pairs are used so that the reported violation rate has an unambiguous denominator. For qualitative plots, the MAE of each test cell is first averaged across the four factorial variants. The representative battery is then selected as the cell whose cross-model mean MAE is closest to the median across test cells, with the battery identifier used as a deterministic tie breaker. Thus, the displayed cell is not selected in favor of either the proposed model or a visually desirable trajectory.
\section{Results and Discussion}
\label{sec:results}

% This result shell intentionally contains no legacy values. Formal values,
% tables, and figures are populated only from non-smoke runs produced after the
% battery-boundary and preprocessing-leakage corrections.
\subsection{Full-supervision reference}
\label{subsec:new_full_supervision}

The $r=1.0$ condition establishes the accuracy ceiling of the four controlled variants when every eligible training target is available. The final table will report test-cell MAE, RMSE, $R^2$, and trajectory metrics as mean $\pm$ sample standard deviation over the predeclared repetition policy. Interpretation will separate two questions: whether the multi-scale representation improves the data-driven CNN-LSTM, and whether the soft monotonic term changes pointwise accuracy or trajectory regularity when labels are dense. A physical-constraint variant will not be described as superior merely because it produces fewer upward deviations; any simultaneous error increase will be reported explicitly.

\subsection{Performance under decreasing SOH-label budgets}
\label{subsec:new_label_budget}

The main result varies $r$ from 1.0 to 0.3 while holding the test cells fully labeled. For each variant and ratio, the result archive supplies one observation per completed seed combination; group means and sample standard deviations are computed only after excluding smoke runs. The principal figure will show MAE and RMSE against $r$, with a separate $R^2$ panel or table. The combined figure and every individual panel are generated independently from the same result summary at 600 dpi and in PDF/SVG format, rather than by cropping a composite image.

The central comparison is degradation relative to each model's own $r=1.0$ reference, supplemented by between-model differences at the same $r$. This distinction prevents a model with a stronger dense-label baseline from appearing label-robust solely because of its absolute error level. The discussion will identify ratios at which the physics term helps, is neutral, or imposes an accuracy cost. Consequently, a label-budget-dependent effect is treated as a substantive finding rather than hidden by averaging all ratios into one score.

\subsection{Trajectory stability on unseen cells}
\label{subsec:new_trajectory_results}

Aggregate point errors do not establish whether a prediction follows a plausible degradation path. This subsection will therefore report monotonicity-violation rate, mean violation excess, cumulative upward excess per test battery, and mean absolute second difference for the same runs used in Section~\ref{subsec:new_label_budget}. These metrics will be interpreted alongside MAE and RMSE, not as substitutes for them. In particular, a smoother curve is beneficial only if it does not erase genuine capacity recovery or introduce a systematic late-life bias.

Representative trajectories are selected by the predeclared cross-model median-cell rule in Section~\ref{subsec:metrics}. For the selected unseen cell, the final figure will show the measured SOH, predictions from the four factorial variants, and an aligned error panel. The selection rule, battery identifier, split seed, mask seed, and training seed will be reported in the figure-source record so that the case study cannot be chosen retrospectively for visual advantage.

\subsection{Factorial ablation and interaction of architecture and physics}
\label{subsec:new_factorial_ablation}

After the main results, the $2\times2$ matrix is used to quantify the independent effects of the multi-scale architecture and soft monotonic constraint. For a lower-is-better error metric $q_{ap}(r)$, where $a\in\{0,1\}$ denotes single- or multi-scale convolution and $p\in\{0,1\}$ denotes the absence or presence of physics, the average architecture and physics improvements are defined as
\begin{align}
\Delta_{\mathrm{MS}}(r) &= \tfrac{1}{2}\left[q_{00}(r)-q_{10}(r)+q_{01}(r)-q_{11}(r)\right],\\
\Delta_{\mathrm{PI}}(r) &= \tfrac{1}{2}\left[q_{00}(r)-q_{01}(r)+q_{10}(r)-q_{11}(r)\right].
\end{align}
Positive values indicate an average reduction in error. The corresponding paired per-run differences will be retained for uncertainty analysis whenever the variants share the same split, mask, and initialization seeds. Trajectory metrics are analyzed in parallel to determine whether an apparent physics benefit arises from improved estimation, reduced nonphysical upward motion, or a trade-off between the two. This ablation is deliberately placed after the principal label-budget and trajectory findings because its role is explanatory rather than evidentiary precedence.

The discussion will close with the evidence boundary of the experiment. HUST random masks support conclusions about controlled label-budget sensitivity and unseen-cell generalization under the tested distribution. They do not by themselves support claims about temporal extrapolation from early-life labels, periodic BMS calibration, cross-chemistry transfer, or field deployment.

\iffalse
\subsection{Full-supervision benchmark on the public dataset}
\label{subsec:full_supervision_results}

The full-supervision experiment first evaluated whether PI-MSCL can provide a reliable SOH estimator under the standard public-dataset setting, where every training sample has an available SOH label. Table~\ref{tab:full_supervision} summarizes the test performance of XGBoost, LSTM, CNN-LSTM, and PI-MSCL under $r=1.0$.

\begin{table}[t]
\caption{Full-supervision benchmark performance on the test cells ($r=1.0$).}
\label{tab:full_supervision}
\centering
\begin{tabular}{lccc}
\toprule
Model & RMSE (\%) & MAE (\%) & $R^2$ \\
\midrule
XGBoost & 1.907 & 1.097 & 0.9322 \\
LSTM & 2.086 & 1.164 & 0.9188 \\
CNN-LSTM & 1.945 & \textbf{1.059} & 0.9294 \\
PI-MSCL & \textbf{1.688} & 1.101 & \textbf{0.9461} \\
\bottomrule
\end{tabular}
\end{table}

XGBoost achieved an RMSE of 1.907\% and an MAE of 1.097\%, indicating that hand-crafted charging features contain useful SOH information. However, its static regression structure does not explicitly model cross-cycle temporal evolution. The LSTM baseline produced higher RMSE and MAE than XGBoost, showing that recurrent modeling alone was not sufficient when local degradation features were not explicitly extracted. CNN-LSTM reduced MAE to 1.059\%, the lowest MAE in this full-supervision comparison, but its RMSE and $R^2$ remained close to those of XGBoost.

PI-MSCL achieved the best RMSE and $R^2$, with RMSE decreasing to 1.688\% and $R^2$ increasing to 0.9461. Its MAE was slightly higher than that of CNN-LSTM, which indicates that the physical consistency constraint did not uniformly reduce all point-wise errors under dense supervision. The gain appeared mainly in suppressing larger trajectory deviations and improving the global consistency of SOH predictions. This observation provides an important baseline for the later incomplete-supervision analysis: PI-MSCL was not designed simply to minimize point-wise MAE under ideal full-label conditions, but to improve trajectory-level robustness when supervision becomes incomplete.

\subsection{Architecture and physical-constraint ablation}
\label{subsec:ablation_results}

To separate the effects of multi-scale representation and physical consistency, a two-dimensional ablation was conducted. The architecture axis compared the single-scale CNN-LSTM with the proposed multi-scale MS-CNN-LSTM. The constraint axis compared models trained without physics and models trained with the soft monotonic physical consistency loss. Table~\ref{tab:ablation_matrix_full} reports the RMSE values under full supervision.

\begin{table}[t]
\caption{Two-dimensional ablation matrix under full supervision. Values are RMSE (\%).}
\label{tab:ablation_matrix_full}
\centering
\begin{tabular}{lccc}
\toprule
Model & No physics & Soft monotonic & Constraint gain \\
\midrule
CNN-LSTM & 1.945 & 1.819 & 0.126 \\
MS-CNN-LSTM & 1.779 & \textbf{1.688} & 0.091 \\
Architecture gain & 0.166 & 0.131 & -- \\
\bottomrule
\end{tabular}
\end{table}

Both architectural improvement and physical consistency reduced RMSE. Replacing the single-scale CNN-LSTM with MS-CNN-LSTM reduced RMSE by 0.166 percentage points without the physical constraint and by 0.131 percentage points with the physical constraint. Adding the soft monotonic constraint reduced RMSE by 0.126 percentage points for the single-scale architecture and by 0.091 percentage points for the multi-scale architecture. These results show that multi-scale representation was the larger contributor under full supervision, while the physical constraint provided a smaller but still positive trajectory-regularization effect.

The same ablation was then extended to all supervision ratios, as shown in Table~\ref{tab:ablation_ratios}. The RMSE improvement from the multi-scale architecture was stable across $r=1.0$, $0.7$, $0.5$, and $0.3$. The full PI-MSCL model consistently achieved the lowest RMSE among the four ablation variants at each ratio.

\begin{table}[t]
\caption{RMSE comparison across supervision ratios in the architecture-constraint ablation.}
\label{tab:ablation_ratios}
\centering
\begin{tabular}{lcccc}
\toprule
Ratio & CNN-LSTM & PI-CNN-LSTM & MS-CNN-LSTM & PI-MSCL \\
\midrule
$r=1.0$ & 1.945 & 1.819 & 1.779 & \textbf{1.688} \\
$r=0.7$ & 2.086 & 1.864 & 1.814 & \textbf{1.728} \\
$r=0.5$ & 2.355 & 2.057 & 1.982 & \textbf{1.876} \\
$r=0.3$ & 2.380 & 2.154 & 2.161 & \textbf{1.937} \\
\bottomrule
\end{tabular}
\end{table}

The ablation results support two conclusions. First, multi-scale convolution was the main source of representational improvement because it reduced RMSE under both physics-free and physics-informed settings. Second, the soft monotonic constraint became more important as the label density decreased. At $r=0.3$, the physical constraint reduced the RMSE of the multi-scale model from 2.161\% to 1.937\%, a larger absolute gain than under full supervision. This behavior is consistent with the role of physical consistency as structural weak supervision when direct SOH labels become sparse.

\subsection{Robustness under partial lifecycle supervision}
\label{subsec:partial_results}

The partial-supervision experiments evaluated how model accuracy changed as the supervision ratio decreased. Table~\ref{tab:mae_ratios} reports MAE values for XGBoost, LSTM, CNN-LSTM, and PI-MSCL under four supervision ratios.

\begin{table}[t]
\caption{MAE (\%) under different supervision ratios.}
\label{tab:mae_ratios}
\centering
\begin{tabular}{lcccc}
\toprule
Model & $r=1.0$ & $r=0.7$ & $r=0.5$ & $r=0.3$ \\
\midrule
XGBoost & 1.097 & 1.139 & 1.174 & 1.238 \\
LSTM & 1.164 & 1.197 & 1.268 & 1.296 \\
CNN-LSTM & \textbf{1.059} & 1.147 & 1.164 & 1.186 \\
PI-MSCL & 1.101 & \textbf{1.115} & \textbf{1.137} & \textbf{1.134} \\
\bottomrule
\end{tabular}
\end{table}

Under full supervision, CNN-LSTM achieved the lowest MAE. When the supervision ratio decreased, the three baselines degraded more clearly. From $r=1.0$ to $r=0.3$, MAE increased by 0.141 percentage points for XGBoost, 0.132 percentage points for LSTM, and 0.127 percentage points for CNN-LSTM. In contrast, PI-MSCL increased by only 0.033 percentage points, from 1.101\% to 1.134\%. The nearly flat MAE curve indicates that PI-MSCL retained stable point-wise accuracy when direct SOH labels were reduced.

PI-MSCL became the best-performing model once the supervision ratio dropped below full supervision. At $r=0.7$, PI-MSCL achieved an MAE of 1.115\%, lower than CNN-LSTM at 1.147\%. At $r=0.5$ and $r=0.3$, PI-MSCL remained the best model, with MAE values of 1.137\% and 1.134\%, respectively. This result supports the central motivation of the framework: the physical constraint is most valuable when the data-fitting term is weakened by sparse labels.

Table~\ref{tab:partial_multimetric} further compares CNN-LSTM and PI-MSCL using MAE, RMSE, and $R^2$ in the partial-supervision regimes. PI-MSCL improved all three metrics at $r=0.7$, $r=0.5$, and $r=0.3$. The RMSE reduction was larger than the MAE reduction. For example, at $r=0.3$, RMSE decreased from 2.211\% for CNN-LSTM to 1.921\% for PI-MSCL, while MAE decreased from 1.186\% to 1.134\%. Since RMSE is more sensitive to large errors, this difference indicates that PI-MSCL mainly reduced larger trajectory deviations rather than uniformly shrinking every point-wise error.

\begin{table*}[t]
\caption{Multi-metric comparison between CNN-LSTM and PI-MSCL under partial supervision.}
\label{tab:partial_multimetric}
\centering
\begin{tabular}{lcccccc}
\toprule
Ratio & \multicolumn{3}{c}{CNN-LSTM} & \multicolumn{3}{c}{PI-MSCL} \\
\cmidrule(lr){2-4}\cmidrule(lr){5-7}
& MAE (\%) & RMSE (\%) & $R^2$ & MAE (\%) & RMSE (\%) & $R^2$ \\
\midrule
$r=0.7$ & 1.147 & 1.989 & 0.9239 & \textbf{1.115} & \textbf{1.707} & \textbf{0.9431} \\
$r=0.5$ & 1.164 & 2.102 & 0.9176 & \textbf{1.137} & \textbf{1.842} & \textbf{0.9283} \\
$r=0.3$ & 1.186 & 2.211 & 0.9045 & \textbf{1.134} & \textbf{1.921} & \textbf{0.9185} \\
\bottomrule
\end{tabular}
\end{table*}

\subsection{Trajectory-level behavior on an unseen cell}
\label{subsec:trajectory_case}

Point-wise metrics do not fully reveal whether an SOH model produces physically reasonable degradation trajectories. A trajectory-level case study was therefore conducted on an unseen test cell under the sparse supervision setting $r=0.3$. The final figure will compare the true SOH trajectory, the CNN-LSTM prediction, and the PI-MSCL prediction for the same cell.

\begin{figure*}[t]
  \centering
  \fbox{\parbox[c][0.18\textheight][c]{0.86\textwidth}{\centering Placeholder for the unseen-cell SOH trajectory comparison.\\ The final figure should show true SOH, CNN-LSTM prediction, and PI-MSCL prediction under $r=0.3$.}}
  \caption{Trajectory-level SOH prediction comparison on an unseen cell under sparse supervision.}
  \label{fig:trajectory_case}
\end{figure*}

In the current manuscript data, the CNN-LSTM trajectory showed stronger local oscillations and a more visible stage-wise bias in the unlabeled lifecycle region. Its RMSE on the representative unseen cell was 0.454\%. PI-MSCL produced a smoother and more monotonic trajectory, with an RMSE of 0.207\% on the same cell. This trajectory-level reduction is consistent with the aggregate RMSE improvements in Table~\ref{tab:partial_multimetric}.

A corresponding error-evolution plot will be inserted after the final figure assets are consolidated. In the current manuscript data, the PI-MSCL error curve stayed closer to zero, with an MAE of 0.159\%, whereas CNN-LSTM showed a larger systematic positive drift in the later lifecycle region, with an MAE of 0.346\%. In SOH estimation, positive drift in late-life regions can be safety-sensitive because it corresponds to overestimating the remaining health of a degrading battery. The case study therefore illustrates the practical value of combining multi-scale feature learning with a soft monotonic physical constraint: the model not only improves aggregate metrics, but also suppresses trajectory shapes that are undesirable for BMS use.
\fi
% Archived from the main evidence chain on 2026-08-06. The underlying draft
% remains recoverable here and in backups/manuscript_pre_restructure_20260806,
% but it is excluded because the synthetic risk experiment is outside the
% controlled sparse-label SOH question evaluated in Sections 3--4.
\iffalse
\section{Trajectory-Based Abnormal Degradation Risk Indication}
\label{sec:risk}

The preceding experiments focused on SOH estimation accuracy and robustness under partial lifecycle supervision. In practical BMS deployment, however, the predicted SOH trajectory may also provide auxiliary information for abnormal degradation risk indication. This section therefore examines whether the trained SOH estimator can be reused, without training an additional fault-diagnosis model, to generate low-cost warning signals for degradation patterns that visibly disturb the predicted trajectory.

The scope of this analysis is deliberately limited. The objective is not to replace electrochemical diagnosis, impedance analysis, or dedicated fault classifiers. Instead, the objective is to test whether a physically constrained SOH prediction trajectory can expose deviations associated with sudden accelerated aging, capacity-knee behavior, and lithium-plating-like step drops. Knee points in battery aging trajectories are widely recognized as important markers of nonlinear degradation transition \cite{attia2022knees}. These scenarios are treated as trajectory-level risk patterns because they alter either the local SOH level, the degradation slope, or both. This framing is consistent with the broader degradation-diagnostics view that battery faults and aging mechanisms often appear as changes in observable capacity and response trajectories \cite{birkl2017degradation}.

\subsection{Trajectory-Deviation Risk Score}

The risk score is constructed from the model output rather than from additional SOH labels. Let $\hat{y}_t$ denote the SOH predicted by a trained estimator at cycle or window $t$. For each detection window, a first-order local trend is fitted to the recent prediction sequence,
\begin{equation}
  \tilde{y}_t = a_t t + b_t,
\end{equation}
where $a_t$ and $b_t$ are obtained by linear regression over the local window. The trajectory-deviation score is then defined as the absolute residual between the current prediction and the local trend,
\begin{equation}
  s_t = |\hat{y}_t - \tilde{y}_t|.
\end{equation}
The underlying assumption is that a normal SOH trajectory should remain locally smooth and approximately monotonic over a short window, whereas an abnormal event or an abrupt degradation-regime change will increase the residual from the local trend.

This score is attractive for deployment because it depends only on the online SOH predictions. No ground-truth SOH label is required at the detection time, and no separate abnormal-state classifier is trained. The score therefore evaluates whether PI-MSCL's smoother and more physically consistent prediction trajectory can increase the signal-to-noise ratio of trajectory disturbances relative to a purely data-driven CNN-LSTM baseline.

\subsection{Thresholding and Warning Rule}

To reduce false alarms caused by isolated noise, the warning rule combines a statistical threshold with a persistence criterion. The detection threshold is estimated from the normal segment of the score sequence as
\begin{equation}
  \tau = \mu_s + 2\sigma_s,
\end{equation}
where $\mu_s$ and $\sigma_s$ denote the mean and standard deviation of the normal trajectory-deviation scores. A warning is triggered only when $s_t > \tau$ holds for $N$ consecutive windows. In the current manuscript experiments, $N$ is set to 20. This N-hit strategy trades a modest detection delay for improved robustness against transient fluctuations.

The detection performance is evaluated using three complementary metrics. The area under the receiver operating characteristic curve (AUC) measures the separability between normal and abnormal score distributions across thresholds. Det@FPR5\% reports the true detection rate when the window-level false-positive rate is fixed at 5\%. Delay reports the number of windows between abnormal injection and the first satisfied N-hit warning condition.

\subsection{Risk-Indication Results}

Figure~\ref{fig:risk_trajectories} will show representative normal and abnormal SOH prediction trajectories for the in-house abnormal-degradation scenarios after the final figure assets are consolidated. In the current manuscript data, sudden accelerated aging produced a clear separation between normal and abnormal predicted trajectories near the injection point, followed by a steeper degradation slope. Capacity-knee behavior showed a milder pre-event trajectory followed by an increased post-knee decay rate. The lithium-plating-like step-drop scenario produced an abrupt SOH decrease followed by continued accelerated decay. These trajectory shapes are consistent with the intended use of the deviation score: the score responds when the predicted trajectory no longer follows its recent local trend.

\begin{figure*}[t]
  \centering
  \fbox{\parbox[c][0.18\textheight][c]{0.86\textwidth}{\centering Placeholder for normal and abnormal SOH prediction trajectories.\\ The final figure should compare CNN-LSTM and PI-MSCL under sudden aging, capacity knee, and lithium-plating-like step-drop scenarios.}}
  \caption{Representative trajectory-level risk patterns under abnormal degradation scenarios.}
  \label{fig:risk_trajectories}
\end{figure*}

Figure~\ref{fig:risk_scores} will present the corresponding trajectory-deviation scores, normal-score baseline, and detection threshold. The current results indicate that PI-MSCL yields larger and more sustained score excursions than CNN-LSTM in abrupt degradation scenarios. This behavior is expected because the physical constraint suppresses background oscillations in normal prediction trajectories, making abnormal deviations easier to separate from normal variation.

\begin{figure*}[t]
  \centering
  \fbox{\parbox[c][0.18\textheight][c]{0.86\textwidth}{\centering Placeholder for trajectory-deviation scores and N-hit warnings.\\ The final figure should show the score baseline, threshold, and warning time for each scenario.}}
  \caption{Trajectory-deviation scores and warning decisions under abnormal degradation scenarios.}
  \label{fig:risk_scores}
\end{figure*}

Table~\ref{tab:risk_detection} summarizes the quantitative warning performance reported in the current manuscript. Across the three abnormal-degradation scenarios and three severity levels, PI-MSCL consistently improves AUC and Det@FPR5\% and reduces detection delay compared with CNN-LSTM. The advantage is particularly clear in abrupt trajectory-disturbance scenarios. For the severe lithium-plating-like step-drop case, PI-MSCL achieves an AUC of 0.893, Det@FPR5\% of 87.3\%, and a mean delay of 5.0 windows, compared with 0.837, 80.8\%, and 18.5 windows for CNN-LSTM. In the medium sudden-aging case, the delay decreases from 11.2 windows to 5.8 windows. In the medium capacity-knee case, Det@FPR5\% increases from 67.2\% to 85.4\%.

\begin{table*}[t]
\caption{Abnormal degradation risk-indication performance under different scenarios.}
\label{tab:risk_detection}
\centering
\begin{tabular}{llcccccc}
\toprule
Scenario & Severity & \multicolumn{3}{c}{CNN-LSTM} & \multicolumn{3}{c}{PI-MSCL} \\
\cmidrule(lr){3-5}\cmidrule(lr){6-8}
 & & AUC & Det@FPR5\% & Delay & AUC & Det@FPR5\% & Delay \\
\midrule
Sudden aging & Mild & 0.716 & 63.5 & 58.8 & 0.809 & 74.3 & 7.2 \\
Sudden aging & Medium & 0.734 & 67.3 & 11.2 & 0.815 & 76.3 & 5.8 \\
Sudden aging & Severe & 0.729 & 71.8 & 8.0 & 0.817 & 78.2 & 4.8 \\
Capacity knee & Mild & 0.720 & 62.6 & 96.0 & 0.837 & 77.9 & 23.0 \\
Capacity knee & Medium & 0.741 & 67.2 & 59.0 & 0.869 & 85.4 & 15.2 \\
Capacity knee & Severe & 0.750 & 68.8 & 67.8 & 0.858 & 83.9 & 21.2 \\
Lithium-plating step drop & Mild & 0.741 & 67.5 & 28.0 & 0.838 & 81.3 & 10.0 \\
Lithium-plating step drop & Medium & 0.756 & 70.8 & 25.0 & 0.847 & 83.4 & 6.0 \\
Lithium-plating step drop & Severe & 0.837 & 80.8 & 18.5 & 0.893 & 87.3 & 5.0 \\
\bottomrule
\end{tabular}
\end{table*}

These results suggest that PI-MSCL is useful not only as a point-wise SOH estimator, but also as a trajectory generator whose physically smoothed output can support simple risk screening. The interpretation should nevertheless remain bounded. The current deviation score is most suitable for abrupt or slope-changing degradation patterns. Progressive anomalies with slowly accumulating deviations may require longer detection windows, multi-signal fusion, or a dedicated diagnostic model. Therefore, the proposed risk-indication module should be viewed as an engineering extension of SOH estimation rather than as a complete battery-fault diagnosis framework.

\fi

\section{Conclusion}
\label{sec:conclusion}

This study formulates lithium-ion battery SOH estimation under a controlled cycle-level label budget and develops PI-MSCL to use information that remains available when a target is masked. Its multi-scale convolutional branches represent degradation patterns over different temporal extents, the LSTM models their cross-cycle evolution, and the soft monotonic term supplies label-independent trajectory structure. The experimental design deliberately separates the multi-scale and physics factors so that an accuracy change cannot be attributed to a different data windowing or preprocessing pipeline.

The final empirical conclusion will be based only on the leakage-free multi-seed runs archived by the unified experiment runner. In particular, the analysis will distinguish pointwise accuracy from trajectory plausibility and will report unfavorable or label-budget-dependent effects of the physical term if they occur. Earlier development summaries are not carried forward because they used inconsistent battery-window construction or preprocessing that accessed unseen-cell statistics.

The scope of the study is bounded in three ways. First, random masking on completely labeled HUST data measures robustness to a controlled label budget; it does not reproduce continuous early-life labeling, periodic BMS calibration, or naturally missing field records. Second, evaluation on one LFP dataset does not establish cross-chemistry or field-deployment generality. Third, soft monotonicity is a structural prior rather than an electrochemical degradation model and may trade pointwise accuracy for trajectory regularity at some label ratios. Future validation should therefore include predeclared prefix and periodic calibration protocols, additional chemistries, and operational datasets with independently established SOH references.

\section*{Data Availability}

The HUST battery dataset is publicly available through Mendeley Data at \url{https://doi.org/10.17632/nsc7hnsg4s.2} \cite{Yuan2022HUSTDataset}. No private or in-house battery dataset is used in the main study.

\section*{Declaration of Competing Interest}

The authors declare that they have no known competing financial interests or personal relationships that could have appeared to influence the work reported in this paper.

\section*{Acknowledgements}

To be completed.

\printcredits

\bibliographystyle{cas-model2-names}
\bibliography{references}

\end{document}








